% ============================================================
% CAPÍTULO 1 — INTRODUÇÃO
% Dissertação: Metodologias Ativas na Educação Brasileira
% Autor: Fernando Ramos Passoni
% ============================================================

\chapter{Introdução}
\label{chap:introducao}

A educação brasileira enfrenta, há décadas, desafios estruturais que comprometem a qualidade do processo de ensino-aprendizagem \cite{BRASIL2018}. Entre esses desafios, destaca-se a predominância de práticas pedagógicas tradicionais, centradas na transmissão unilateral de conhecimento, que pouco estimulam o pensamento crítico, a resolução de problemas e a capacidade de os discentes aplicarem os conteúdos em contextos reais \cite{BACICH2018}. Esse cenário evidencia a necessidade urgente de repensar os paradigmas educacionais vigentes, sobretudo em um momento em que as demandas sociais, tecnológicas e profissionais se transformam a uma velocidade sem precedentes, exigindo competências que vão além da mera memorização de conteúdos disciplinares.

Nesse contexto, as metodologias ativas — entre as quais se destacam o Aprendizado Baseado em Problemas (ABP) e o Aprendizado Baseado em Projetos (ABPr) — surgem como alternativas pedagógicas promissoras, capazes de reconfigurar a dinâmica das salas de aula ao colocar o estudante no centro do processo educacional. Essas abordagens, amplamente pesquisadas e aplicadas em diversos países, têm demonstrado potencial para estimular a aprendizagem significativa, o trabalho colaborativo e o desenvolvimento de competências transferíveis para a vida profissional e cidadã \cite{BACICH2018}. No entanto, sua implementação no contexto brasileiro ainda enfrenta resistências de ordem institucional, cultural e financeira, o que torna necessária uma investigação aprofundada sobre os fatores que favorecem ou dificultam sua adoção em diferentes realidades educacionais do país.

A presente dissertação insere-se no Programa de Pós-Graduação em Tecnologia Educacional (PPGTE) da Universidade Federal do Ceará (UFC), buscando contribuir para o avanço do conhecimento na interface entre inovação pedagógica e políticas educacionais. A escolha por investigar o ABP e o ABPr justifica-se pelo fato de que essas metodologias representam modelos contrastantes e, ao mesmo tempo, complementares de aprendizagem ativa, cuja análise comparativa pode revelar sinergias e especificidades relevantes para a formulação de propostas educacionais mais robustas e contextualizadas.

\section{Objetivo e finalidade do projeto}
\label{sec:objetivo_finalidade}

O presente projeto de dissertação tem por finalidade investigar, de modo sistemático e aprofundado, o estado da arte, os desafios e os impactos das metodologias ativas — em especial o Aprendizado Baseado em Problemas (ABP) e o Aprendizado Baseado em Projetos (ABPr) — no contexto da educação brasileira, com vistas a formular propostas concretas para uma transformação equitativa dos processos educativos (ver \ref{chap:aspectos_estrategicos}). A investigação proposta insere-se no escopo do PPGTE/UFC e busca responder a uma lacuna identificada na literatura brasileira: a escassez de estudos comparativos que articulem os fundamentos teóricos dessas metodologias com evidências empíricas de sua aplicação em instituições de ensino superior nacionais. A finalidade última do projeto é produzir conhecimento acadêmico de qualidade que possa orientar gestores educacionais, docentes e formuladores de políticas públicas na adoção de práticas pedagógicas mais eficazes e inclusivas.

% ----------------------------------------------------------------
\subsection{Problemas sociais do tema no contexto brasileiro}
\label{subsec:problemas_sociais}

\begin{itemize}
    \item Desigualdade de acesso e permanência na educação básica e superior;
    \item Baixo desempenho em avaliações nacionais e internacionais (PISA, ENEM);
    \item Defasagem entre a formação oferecida pelas instituições de ensino e as demandas do mercado de trabalho;
    \item Resistência institucional à adoção de inovações pedagógicas;
    \item Falta de formação continuada docente voltada para metodologias ativas;
    \item Inexistência de uma cultura institucional que valorize a experimentação pedagógica e a reflexão crítica sobre práticas educativas;
    \item Fragilidade dos mecanismos de avaliação da aprendizagem, frequentemente pautados em provas escritas de caráter memorístico;
    \item Inadequação da infraestrutura física e tecnológica de muitas instituições públicas para suportar metodologias colaborativas e mediadas por tecnologia.
\end{itemize}

Esses problemas, interligados e de natureza sistêmica, configuram um quadro em que a inovação pedagógica se depara com obstáculos que transcendem a esfera individual dos docentes e dos estudantes, alcançando dimensões institucionais, políticas e culturais. A compreensão dessas barreiras é fundamental para que as propostas de implementação de metodologias ativas considerem não apenas os aspectos técnicos e didáticos, mas também as condições estruturais e os fatores socioeconômicos que condicionam a prática educacional brasileira.

% ----------------------------------------------------------------
\subsection{Solução geral}
\label{subsec:solucao_geral}

A solução geral consiste na adoção de metodologias ativas como ABP e ABPr como estratégias pedagógicas estruturantes, complementadas por práticas de avaliação formativa e pelo uso de tecnologias educacionais, visando melhorar significativamente a qualidade da aprendizagem e a equidade no sistema educacional brasileiro. Essa abordagem pressupõe uma reconfiguração das relações pedagógicas, na qual o docente assume o papel de mediador e facilitador do conhecimento, enquanto o discente é estimulado a assumir responsabilidade pelo próprio processo de aprendizagem, desenvolvendo competências cognitivas, socioemocionais e profissionais de forma integrada.

A implementação dessas metodologias requer, contudo, um conjunto de condições que incluem formação docente qualificada, infraestrutura adequada, flexibilidade curricular e comprometimento institucional. A solução geral proposta neste projeto prevê a articulação desses elementos por meio de um modelo de referência que possa ser adaptado a diferentes contextos educacionais, desde escolas da educação básica até cursos de graduação e pós-graduação, considerando as especificidades regionais e institucionais do cenário brasileiro.

% ----------------------------------------------------------------
\subsection{Solução particular}
\label{subsec:solucao_particular}

A solução particular deste estudo compreende a elaboração de um modelo de implementação e avaliação de ABP e ABPr, validado por meio de estudo de caso em instituições de ensino superior brasileiras, que contempla aspectos como planejamento curricular, formação docente, infraestrutura necessária e indicadores de desempenho (ver \ref{chap:organizacao_trabalho}). O recorte metodológico adotado privilegia a perspectiva qualitativa, complementada por dados quantitativos, permitindo uma compreensão densa dos fenômenos educacionais investigados, sem abrir mão da possibilidade de generalização controlada dos achados.

Os contextos de aplicação compreendem instituições públicas e privadas de diferentes regiões do país, selecionadas de acordo com critérios de diversidade institucional e de grau de implementação de metodologias ativas. As ferramentas de análise incluem entrevistas semiestruturadas com docentes e gestores, observação participante em salas de aula, análise documental de projetos pedagógicos e aplicação de questionários com discentes, todos organizados em torno de categorias teóricas previamente definidas a partir da revisão de literatura.

% ----------------------------------------------------------------

\section{Objetivos do projeto}
\label{sec:objetivos_projeto}

A definição dos objetivos do projeto segue uma lógica que vai do geral ao específico, articulando finalidades amplas de natureza teórica e prática com metas concretas de investigação. Essa estruturação visa assegurar que a pesquisa contribua simultaneamente para o avanço do conhecimento acadêmico e para a transformação das práticas educacionais nas instituições investigadas.

\subsection{Objetivos específicos do projeto}
\label{subsec:obj_especificos_projeto}

\begin{enumerate}
    \item Mapear o estado da arte das metodologias ativas (ABP e ABPr) na educação brasileira;
    \item Identificar os principais desafios e barreiras para a implementação dessas metodologias;
    \item Analisar os impactos da adoção de ABP e ABPr sobre a aprendizagem dos discentes;
    \item Formular propostas concretas para a implementação equitativa dessas metodologias.
\end{enumerate}

% ----------------------------------------------------------------
\subsection{Objetivos específicos da pesquisa}
\label{subsec:obj_especificos_pesquisa}

No âmbito da pesquisa propriamente dita, os objetivos específicos contemplam as etapas metodológicas necessárias para o cumprimento das metas do projeto, abrangendo desde a revisão sistemática da literatura até a coleta e análise de dados em campo, passando pela construção de instrumentos de avaliação e pela elaboração de recomendações para políticas educacionais.

\begin{enumerate}
    \item Realizar revisão bibliográfica e estado da arte sobre ABP e ABPr no Brasil e no exterior;
    \item Coletar e analisar dados qualitativos e quantitativos em instituições de ensino;
    \item Avaliar a eficácia das metodologias ativas por meio de indicadores pré e pós-intervenção;
    \item Elaborar um relatório integrado com recomendações para políticas educacionais.
\end{enumerate}

% ----------------------------------------------------------------

\section{Conteúdo e alcance do projeto}
\label{sec:conteudo_alcance}

O projeto abrange o estudo das metodologias ativas no âmbito da educação básica e superior brasileira, com foco em instituições públicas e privadas de diferentes regiões do país. O período de análise compreende os últimos cinco anos, permitindo a identificação de tendências contemporâneas e a avaliação de resultados de intervenções recentes. Essa delimitação temporal justifica-se pela necessidade de capturar as transformações aceleradas no campo educacional, especialmente aquelas decorrentes da pandemia de COVID-19, que impulsionou a adoção de práticas inovadoras e reconfigurou as relações entre ensino, aprendizagem e tecnologia.

Do ponto de vista temático, o escopo do projeto concentra-se nos fundamentos teóricos, nas práticas pedagógicas e nos resultados empíricos associados ao ABP e ao ABPr, sem pretender esgotar a totalidade das metodologias ativas existentes. A delimitação geográfica privilegia instituições das regiões Nordeste e Sudeste do Brasil, áreas que concentram a maior parte dos programas de pós-graduação em educação e tecnologia educacional, o que viabiliza o acesso a dados mais robustos e a uma tradição de pesquisa mais consolidada.

\subsection{Limitações e obstáculos}
\label{subsec:limitacoes}

\begin{itemize}
    \item Dificuldade de acesso a dados institucionais detalhados e padronizados;
    \item Resistência de alguns atores educacionais à mudança de práticas consolidadas;
    \item Limitações orçamentárias para coleta de dados em múltiplas localidades;
    \item Necessidade de aprovação por comitês de ética em pesquisa;
    \item Variabilidade dos contextos institucionais, que pode dificultar a generalização dos resultados.
\end{itemize}

Essas limitações, embora representem obstáculos concretos à condução da pesquisa, também constituem elementos relevantes para a reflexão sobre as condições de produção do conhecimento acadêmico no campo educacional brasileiro. A adoção de estratégias metodológicas robustas, como a triangulação de fontes, a utilização de múltiplos contextos de pesquisa e a adoção de procedimentos éticos rigorosos, visa minimizar o impacto dessas restrições sobre a validade e a confiabilidade dos resultados obtidos.

% ----------------------------------------------------------------

\section{Resultados esperados}
\label{sec:resultados_esperados}

\begin{itemize}
    \item Modelo de implementação e avaliação de ABP e ABPr para instituições de ensino brasileiras;
    \item Publicação de artigo(s) em periódicos Qualis A1/A2;
    \item Contribuição teórica para o campo da educação e da pedagogia;
    \item Recomendações para políticas públicas de educação;
    \item Instrumentos de avaliação replicáveis por outras instituições.
\end{itemize}

Esses resultados esperados articulam dimensões teóricas e práticas, refletindo a vocação do programa de pós-graduação em produzir conhecimento acadêmico com impacto social. A expectativa é que o modelo proposto possa ser utilizado como referência por instituições que estejam em processo de implementação ou reformulação de práticas pedagógicas ativas, contribuindo para a difusão de uma cultura educacional mais inclusiva, colaborativa e orientada para o desenvolvimento de competências.

% ----------------------------------------------------------------

\section{Descrição do caso}
\label{sec:descricao_caso}

\subsection{Aprendizado baseado em problemas (ABP)}
\label{subsec:abp_descricao}

O Aprendizado Baseado em Problemas (ABP) é uma abordagem pedagógica que tem suas raízes nas escolas médicas canadenses da década de 1960, mais especificamente na McMaster University \cite{BARROWS1996}. Essa metodologia parte de problemas autênticos, contextualizados e relevantes, que servem como ponto de partida para a construção do conhecimento pelo estudante, que aprende de forma ativa, colaborativa e reflexiva \cite{SCHMIDT2001}. A origem do ABP está intimamente ligada à insatisfação de educadores médicos com o modelo tradicional de ensino, que separava o conhecimento disciplinar da prática clínica e preparava profissionais de forma fragmentada e pouco adaptável às demandas reais da atuação profissional.

Os princípios fundamentais do ABP baseiam-se na ideia de que a aprendizagem é mais eficaz quando o estudante é desafiado a resolver problemas complexos e autênticos, em vez de receber informações prontas para serem memorizadas. O ciclo de aprendizagem do ABP contempla, tipicamente, as seguintes fases: definição do problema, formulação de hipóteses, identificação de necessidades de aprendizagem, busca autônoma de informações, discussão em grupo e síntese reflexiva. Esse ciclo é reapresentado iterativamente, de modo que o estudante desenvolva não apenas conhecimentos disciplinares, mas também habilidades de raciocínio crítico, trabalho em equipe e comunicação \cite{BARROWS1996}.

No contexto brasileiro, o ABP tem sido adotado em diversas instituições de ensino superior, sobretudo nos cursos de ciências da saúde, como medicina, enfermagem e odontologia, onde sua eficácia tem sido amplamente documentada \cite{SCHMIDT2001}. No entanto, sua aplicação tem se expandido gradualmente para outras áreas do conhecimento, como engenharias, ciências sociais e educação, demonstrando a versatilidade e o potencial de adaptação dessa metodologia a diferentes campos disciplinares e contextos institucionais. Essa expansão, embora promissora, ainda enfrenta desafios relacionados à formação docente, à resistência cultural e à inadequação de modelos de avaliação tradicionais para capturar os tipos de aprendizagem promovidos pelo ABP.

% ----------------------------------------------------------------
\subsection{Aprendizado baseado em projetos (ABPr)}
\label{subsec:abpr_descricao}

O Aprendizado Baseado em Projetos (ABPr) é uma metodologia ativa em que os estudantes desenvolvem projetos complexos, integrando conhecimentos de múltiplas áreas, para resolver problemas reais ou simular situações autênticas \cite{THOMAS2000}. Essa abordagem valoriza a interdisciplinaridade, a colaboração e a autonomia dos discentes, promovendo uma aprendizagem significativa e contextualizada \cite{BARRON1998}. As origens do ABPr podem ser rastreadas nas teorias da aprendizagem de John Dewey, que defendia a educação pela experiência e a importância de conectar o ensino aos interesses e às necessidades reais dos estudantes.

Os princípios orientadores do ABPr incluem a centralidade do estudante no processo de aprendizagem, a intencionalidade pedagógica do projeto, a autenticidade da tarefa, a necessidade de reflexão crítica e a produção de um produto ou resultado tangível. As etapas do ABPr geralmente compreendem a escolha do tema, a formulação de questões norteadoras, o planejamento das atividades, a execução do projeto, a apresentação dos resultados e a avaliação reflexiva do processo. Essa estrutura permite que os estudantes desenvolvam competências como pensamento crítico, resolução de problemas, comunicação eficaz e capacidade de trabalhar colaborativamente em contextos diversificados.

No cenário educacional brasileiro, o ABPr tem sido utilizado tanto na educação básica quanto no ensino superior, com destaque para iniciativas em escolas públicas e privadas que buscam integrar conteúdos curriculares com demandas sociais e comunitárias. A metodologia tem se mostrado especialmente relevante para a formação de cidadãos críticos e engajados, na medida em que estimula a reflexão sobre questões sociais, ambientais e culturais por meio da ação concreta. Contudo, a implementação do ABPr no Brasil também enfrenta obstáculos significativos, como a sobrecarga curricular, a falta de tempo docente para o planejamento colaborativo e a dificuldade de avaliar processos de aprendizagem que extrapolam os limites das avaliações tradicionais \cite{THOMAS2000}.

% ============================================================
% Fim do Capítulo 1 — Introdução
% ============================================================
